\subsection{One coordinated referral after offers}\label{sec:two-referral}

Suppose the platform now observes an assignment $M_0$ stable on the initial lists. Join each unmatched applicant to each mutually acceptable program with a current vacancy. These edges are hidden: if one were displayed it would block $M_0$. Give each program its remaining capacity and let $r$ be the size of a maximum matching in this graph. With lists of $\min(k,d_i)$ programs, where $k\ge1$ and $d_i$ is feasible degree, a hidden option can always replace a displayed one. If an entitlement is protected outside the quota, the substitution is made in a discretionary slot. In both replays, every selected applicant had a full five-slot discretionary list.

\begin{theorem}\label{thm:two-referral}
Among extensions that keep every initial contract in a single stable assignment, the maximum number of additional recipients is $r$. One shortlist substitution per selected applicant attains it. Every stable outcome on the revised display gives each consenting applicant their named referral, even if others decline. Under applicant-proposing DA, every original recipient also receives a weakly preferred assignment.
\end{theorem}
\begin{proof}
The new contracts in any extension must form a residual matching. Conversely, choose any such matching $A$. For each consenting participant substitute their assigned program for their last discretionary displayed option. Write $W$ for the original assignment together with the accepted referrals. All new pairs in the display belong to $W$. On old edges the applicant weakly improved. An originally full program has the same holders. At an originally slack one, every applicant with an old edge weakly preferred their original assignment, by stability. Hence $W$ is stable on this display. DA weakly dominates $W$ for applicants. All stable selections have the same matched people \citep{roth1986}. If some stable selection omitted an accepted referral $e$, it would keep its recipient matched even after deleting that unused edge. On that same graph $W\setminus\{e\}$ is stable by the preceding argument, since this applicant now has a subset of their original list. That would leave the person unmatched in one stable outcome only, a contradiction.
\end{proof}
The maximum fixes the original contracts in a stable witness. Opening new options to incumbents would be a different intervention. The empirical calculation also requires deciding which hidden options a person would actually accept. We report both a top-five rank limit and all mutually acceptable vacancies. Work on facilitating matches, adding school seats and guaranteed school referrals provides related interventions \citep{facilitate,tworound,lepesure}.
